\documentclass[11pt, a4paper]{article}

% ─── Encodage et langue ───
\usepackage[utf8]{inputenc}
\usepackage[T1]{fontenc}
\usepackage[french]{babel}

% ─── Maths ───
\usepackage{amsmath, amssymb, amsthm}

% ─── Mise en page ───
\usepackage[margin=2.5cm]{geometry}
\usepackage{enumitem}
\usepackage{booktabs}
\usepackage{float}
\usepackage{xcolor}
\usepackage{tikz}
\usetikzlibrary{arrows.meta, positioning, fit, backgrounds, calc}
\usepackage{hyperref}
\hypersetup{colorlinks=true, linkcolor=blue!60!black, urlcolor=blue!70!black}
% ─── Commandes ───
\newcommand{\R}{\mathbb{R}}
\newcommand{\Ltwo}{L^2(\mathcal{T})}
\newcommand{\Dp}{D_p}
\newcommand{\Dw}{D_w}
\newcommand{\DK}{D_K}

\theoremstyle{definition}
\newtheorem{definition}{Définition}[section]
\theoremstyle{remark}
\newtheorem{remarque}{Remarque}[section]

% ─── Couleurs thème ───
\definecolor{func}{HTML}{2166AC}
\definecolor{vect}{HTML}{B2182B}
\definecolor{mix}{HTML}{4DAF4A}
\definecolor{gris}{HTML}{666666}

% ============================================================================
\title{%
  \textbf{Architecture des distances} \\[0.3em]
  \Large Pour la classification non supervisée de données mixtes \\[0.2em]
  \large Fonctionnel $\times$ Vectoriel
}
\author{Pierre}
\date{Mars 2026}

\begin{document}
\maketitle

\begin{abstract}
Ce document centralise l'étude de toutes les distances utilisées dans le
pipeline de classification non supervisée de données mixtes
(fonctionnelles + vectorielles). On y définit chaque distance, ses paramètres,
son interprétation géométrique, et comment elles s'articulent dans une
architecture à deux niveaux. Le document est auto-suffisant et indépendant
du jeu de données.
\end{abstract}

\tableofcontents
\newpage

% ============================================================================
% ============================================================================
\part{Le problème : mesurer la similarité entre objets mixtes}
% ============================================================================
% ============================================================================

\section{Nature des données}

Chaque individu $i = 1, \ldots, n$ possède deux types d'information :
\[
  \text{Individu } i : \quad
  \underbrace{\hat{X}_i(t) \in \Ltwo}_{\text{courbe (fonctionnel)}}
  \quad + \quad
  \underbrace{Z_i = (z_{i1}, \ldots, z_{ip}) \in \R^p}_{\text{vecteur (classique)}}.
\]

\textbf{Le défi.} Ces deux objets vivent dans des \textbf{espaces différents}
($\Ltwo$ vs $\R^p$). On ne peut pas les empiler dans un même vecteur ni
calculer une distance euclidienne entre eux. Il faut une architecture de
distances qui :
\begin{enumerate}[nosep]
  \item mesure la dissimilarité \textbf{fonctionnelle} (entre courbes),
  \item mesure la dissimilarité \textbf{vectorielle} (entre vecteurs de $\R^p$),
  \item \textbf{combine} les deux de manière contrôlée.
\end{enumerate}

\section{Vue d'ensemble}

L'architecture est organisée en \textbf{deux niveaux} hiérarchiques :

\begin{figure}[H]
\centering
\begin{tikzpicture}[
    base/.style={draw, rounded corners=4pt, minimum width=2.8cm, minimum height=1cm, align=center, font=\small\bfseries},
    fleche/.style={-{Stealth[length=3mm]}, thick},
    etiq/.style={font=\footnotesize\itshape, text=gris},
  ]

  % Briques de base
  \node[base, fill=func!15, text=func] (D0) {$D_0$\\{\footnotesize niveau}};
  \node[base, fill=func!15, text=func, right=1.2cm of D0] (D1) {$D_1$\\{\footnotesize forme}};
  \node[base, fill=vect!15, text=vect, right=3cm of D1] (Ds) {$D_s$\\{\footnotesize vectoriel}};

  % Dp
  \node[base, fill=func!25, text=func, below=1.5cm of $(D0)!0.5!(D1)$] (Dp) {$\Dp(\alpha)$\\{\footnotesize combinée}};

  % Dw et DK
  \node[base, fill=mix!20, text=mix!70!black, below=1.5cm of $(Dp)!0.5!(Ds)$, xshift=-1.2cm] (Dw) {$\Dw(\alpha, \omega)$\\{\footnotesize pondérée}};
  \node[base, fill=mix!20, text=mix!70!black, right=1.5cm of Dw] (DK) {$\DK(\alpha)$\\{\footnotesize noyaux}};

  % Flèches
  \draw[fleche, func!70] (D0) -- (Dp) node[etiq, midway, left=2pt] {$1-\alpha$};
  \draw[fleche, func!70] (D1) -- (Dp) node[etiq, midway, right=2pt] {$\alpha$};
  \draw[fleche, func!70] (Dp) -- (Dw) node[etiq, midway, left=2pt] {$\omega$};
  \draw[fleche, vect!70] (Ds) -- (Dw) node[etiq, midway, right=2pt] {$1-\omega$};
  \draw[fleche, func!70] (Dp) -- (DK) node[etiq, midway, left=2pt] {$K_f$};
  \draw[fleche, vect!70] (Ds) -- (DK) node[etiq, midway, right=2pt] {$K_s$};

  % Niveaux
  \node[font=\footnotesize\scshape, text=gris, left=0.3cm of D0, yshift=0cm] {Niveau 0};
  \node[font=\footnotesize\scshape, text=gris, left=0.3cm of Dp, yshift=0cm] {Niveau 1};
  \node[font=\footnotesize\scshape, text=gris, left=0.3cm of Dw, yshift=0cm] {Niveau 2};

  % Cadres de fond
  \begin{scope}[on background layer]
    \node[fit=(D0)(D1), draw=func!30, dashed, rounded corners=8pt, inner sep=10pt, fill=func!3] {};
    \node[fit=(Ds), draw=vect!30, dashed, rounded corners=8pt, inner sep=10pt, fill=vect!3] {};
  \end{scope}

\end{tikzpicture}
\caption{Architecture à deux niveaux. \textcolor{func}{Bleu} = fonctionnel,
\textcolor{vect}{rouge} = vectoriel, \textcolor{mix}{vert} = mixte.}
\label{fig:archi}
\end{figure}

\newpage
% ============================================================================
% ============================================================================
\part{Niveau 0 — Les briques de base}
% ============================================================================
% ============================================================================

% ────────────────────────────────────────────────────────────────────────
\section{$D_0$ : distance $L^2$ entre courbes}
% ────────────────────────────────────────────────────────────────────────

\begin{definition}[$D_0$]
\[
  D_0(i,j) = \left\| \hat{X}_i - \hat{X}_j \right\|_{L^2}
  = \sqrt{\int_{\mathcal{T}} \bigl[\hat{X}_i(t) - \hat{X}_j(t)\bigr]^2 \, dt}
\]
\end{definition}

\paragraph{Ce qu'elle mesure.}
La \textbf{proximité globale des profils}. C'est la racine carrée de l'aire
entre les deux courbes (au carré). Deux courbes au même \textbf{niveau}
mais légèrement décalées en phase seront proches en $D_0$.

\paragraph{Forces.}
\begin{itemize}[nosep]
  \item Interprétation intuitive : <<~l'aire entre les courbes~>>.
  \item Pas de paramètre.
  \item Exploite la norme naturelle de $\Ltwo$.
\end{itemize}

\paragraph{Faiblesses.}
\begin{itemize}[nosep]
  \item Sensible uniquement au \textbf{niveau} (écart vertical).
  \item Deux courbes de même forme mais translatées verticalement seront
    éloignées, même si leur dynamique est identique.
  \item Insensible aux différences de \textbf{forme} si les niveaux sont
    proches.
\end{itemize}

\paragraph{Calcul numérique.}
L'intégrale est approchée par somme de Riemann sur $M = 1000$ points
équidistants :
\[
  D_0(i,j) \approx \sqrt{\sum_{m=1}^{M} \bigl[\hat{X}_i(t_m) - \hat{X}_j(t_m)\bigr]^2 \cdot \Delta t}, \qquad
  \Delta t = \frac{|\mathcal{T}|}{M - 1}.
\]

% ────────────────────────────────────────────────────────────────────────
\section{$D_1$ : distance $L^2$ entre dérivées}
% ────────────────────────────────────────────────────────────────────────

\begin{definition}[$D_1$]
\[
  D_1(i,j) = \left\| \hat{X}_i' - \hat{X}_j' \right\|_{L^2}
  = \sqrt{\int_{\mathcal{T}} \bigl[\hat{X}_i'(t) - \hat{X}_j'(t)\bigr]^2 \, dt}
\]
\end{definition}

\paragraph{Ce qu'elle mesure.}
La proximité des \textbf{dynamiques}. Deux courbes au même niveau mais dont
l'une accélère plus vite seront proches en $D_0$ mais éloignées en $D_1$.

\paragraph{Interprétation physique.}
Si $X_i(t)$ est une température : $X_i'(t)$ est la \textbf{vitesse de
changement} de la température. Deux stations dont la température chute à des
rythmes très différents en automne auront un $D_1$ élevé, même si leurs
températures moyennes sont similaires.

\paragraph{Forces.}
\begin{itemize}[nosep]
  \item Capture la \textbf{forme} indépendamment du niveau.
  \item Utile quand les dynamiques importent plus que les valeurs absolues.
\end{itemize}

\paragraph{Faiblesses.}
\begin{itemize}[nosep]
  \item Sensible au \textbf{bruit} : la dérivée amplifie les fluctuations
    rapides. Un lissage insuffisant rend $D_1$ très bruitée.
  \item Insensible aux différences de niveau : deux courbes parallèles
    (même forme, niveau différent) ont $D_1 = 0$.
  \item Nécessite un lissage préalable (les données brutes ne sont pas
    dérivables).
\end{itemize}

\begin{remarque}[Lien avec la pénalité de lissage]
La pénalité $\lambda \int [\hat{X}''(t)]^2 dt$ du lissage (étape 01)
contrôle la régularité de la dérivée \emph{seconde}. Si $\lambda$ est trop
petit, les dérivées \emph{premières} sont erratiques et $D_1$ est dominée
par le bruit. Le choix de $\lambda$ par GCV est donc crucial pour la qualité
de $D_1$.
\end{remarque}

\paragraph{Échelle.}
$D_0$ et $D_1$ sont sur des échelles \textbf{très différentes}. Sur Canadian
Weather : $D_0 \in [3.9, 532]$ mais $D_1 \in [0.51, 6.61]$ (facteur
$\sim\!80$). D'où la nécessité de normaliser avant toute combinaison.

% ────────────────────────────────────────────────────────────────────────
\section{$D_s$ : distance euclidienne sur $Z$}
% ────────────────────────────────────────────────────────────────────────

\begin{definition}[$D_s$]
\[
  D_s(i,j) = \left\| Z_i^{\text{std}} - Z_j^{\text{std}} \right\|_2
  = \sqrt{\sum_{l=1}^{p} \left( \frac{z_{il} - \bar{z}_l}{\sigma_{z_l}}
    - \frac{z_{jl} - \bar{z}_l}{\sigma_{z_l}} \right)^2}
\]
où $Z^{\text{std}} = (Z - \bar{Z}) / \sigma_Z$ est la standardisation
variable par variable.
\end{definition}

\paragraph{Ce qu'elle mesure.}
La proximité dans l'espace des variables classiques. Pour Canadian Weather :
latitude, longitude, précipitations moyennes.

\paragraph{Pourquoi standardiser~?}
Sans standardisation, la latitude (en degrés, $\sim\!40\text{--}80$) dominerait
les précipitations (en mm, $\sim\!1\text{--}4$). La standardisation met toutes
les variables sur une même échelle (moyenne 0, variance 1).

\paragraph{Forces.}
\begin{itemize}[nosep]
  \item Simple, pas de paramètre.
  \item Interprétable : distance euclidienne classique.
  \item Souvent très discriminante quand les variables $Z$ sont pertinentes.
\end{itemize}

\paragraph{Faiblesses.}
\begin{itemize}[nosep]
  \item Ignore totalement les courbes.
  \item Suppose que toutes les variables $Z$ sont d'importance égale (après
    standardisation), ce qui n'est pas garanti.
\end{itemize}

\newpage
% ============================================================================
% ============================================================================
\part{Niveau 1 — La distance fonctionnelle combinée $\Dp(\alpha)$}
% ============================================================================
% ============================================================================

\section{Motivation}

$D_0$ capture le \textbf{niveau}, $D_1$ capture la \textbf{forme}. Pourquoi
choisir ? On peut faire les deux avec un paramètre~$\alpha$.

\section{Définition}

\begin{definition}[$\Dp(\alpha)$]
Soient $\tilde{D}_0 = D_0 / \max(D_0)$ et $\tilde{D}_1 = D_1 / \max(D_1)$
les distances normalisées par leur maximum. Alors :
\[
  \boxed{
    \Dp(\alpha) = \sqrt{(1 - \alpha) \cdot \tilde{D}_0^{\,2}
    + \alpha \cdot \tilde{D}_1^{\,2}},
    \qquad \alpha \in [0, 1]
  }
\]
\end{definition}

\section{Interprétation géométrique}

$\Dp(\alpha)$ est une \textbf{norme pondérée} dans l'espace produit
$(\tilde{D}_0, \tilde{D}_1)$. Quand on trace les iso-distances de $\Dp$ dans
le plan $(\tilde{D}_0, \tilde{D}_1)$, on obtient des \textbf{ellipses} dont l'axe
principal tourne avec $\alpha$ :

\begin{center}
\begin{tabular}{lcl}
  \toprule
  $\boldsymbol{\alpha}$ & \textbf{Poids} & \textbf{Signification} \\
  \midrule
  $0$ & $(1, 0)$ & Seul le \textbf{niveau} compte $\Rightarrow \Dp = \tilde{D}_0$ \\
  $0.3$ & $(0.7, 0.3)$ & Surtout le niveau, un peu la forme \\
  $0.5$ & $(0.5, 0.5)$ & Compromis \textbf{50/50} \\
  $0.7$ & $(0.3, 0.7)$ & Surtout la forme \\
  $1$ & $(0, 1)$ & Seule la \textbf{forme} compte $\Rightarrow \Dp = \tilde{D}_1$ \\
  \bottomrule
\end{tabular}
\end{center}

\section{Pourquoi normaliser par le maximum~?}

$D_0 \in [3.9, 532]$ et $D_1 \in [0.51, 6.61]$. Sans normalisation, avec
$\alpha = 0.5$ :
\[
  \Dp \approx \sqrt{0.5 \cdot 532^2 + 0.5 \cdot 6.6^2} \approx 376,
\]
— $D_1$ ne pèse rien. Après normalisation, $\tilde{D}_0, \tilde{D}_1 \in [0,1]$
et $\alpha = 0.5$ donne un vrai compromis 50/50.

\begin{remarque}[Alternative : normalisation par la médiane]
On pourrait normaliser par la médiane au lieu du maximum. Le max est sensible
aux outliers, mais il garantit que $\tilde{D} \in [0, 1]$, ce qui rend
l'interprétation de $\alpha$ plus transparente.
\end{remarque}

\section{$\Dp$ est une distance}

$\Dp(\alpha)$ vérifie les axiomes d'une distance (pour tout $\alpha \in [0,1]$) :
\begin{enumerate}[nosep]
  \item \textbf{Positivité :} $\Dp \geq 0$, et $\Dp(i,j) = 0 \iff i = j$.
  \item \textbf{Symétrie :} $\Dp(i,j) = \Dp(j,i)$.
  \item \textbf{Inégalité triangulaire :} découle de l'inégalité de Minkowski
    appliquée à $\ell^2$ pondéré.
\end{enumerate}
C'est important car PAM nécessite une \textbf{vraie} matrice de distances.

\newpage
% ============================================================================
% ============================================================================
\part{Niveau 2 — Les distances mixtes}
% ============================================================================
% ============================================================================

On a maintenant $\Dp(\alpha)$ (fonctionnel) et $D_s$ (vectoriel). Deux
stratégies pour les combiner :

% ────────────────────────────────────────────────────────────────────────
\section{$\Dw(\alpha, \omega)$ : pondération directe (stratégie B)}
% ────────────────────────────────────────────────────────────────────────

\subsection{Définition}

\begin{definition}[$\Dw$]
Soit $\tilde{D}_s = D_s / \max(D_s)$. Alors :
\[
  \boxed{
    \Dw(\alpha, \omega) = \sqrt{\omega \cdot \Dp(\alpha)^2
    + (1 - \omega) \cdot \tilde{D}_s^{\,2}},
    \qquad \alpha, \omega \in [0, 1]
  }
\]
\end{definition}

\subsection{Les deux boutons}

\begin{center}
\begin{tikzpicture}[
  bouton/.style={draw, circle, minimum size=1.5cm, thick, font=\large\bfseries},
]
  \node[bouton, fill=func!15, text=func] (alpha) {$\alpha$};
  \node[right=0.3cm of alpha, font=\small, text=gris, align=left] {
    Niveau $\longleftrightarrow$ Forme\\
    {\footnotesize (dans le fonctionnel)}
  };

  \node[bouton, fill=mix!15, text=mix!70!black, right=4cm of alpha] (omega) {$\omega$};
  \node[right=0.3cm of omega, font=\small, text=gris, align=left] {
    Fonctionnel $\longleftrightarrow$ Vectoriel\\
    {\footnotesize (entre les deux sources)}
  };
\end{tikzpicture}
\end{center}

\begin{center}
\begin{tabular}{cll}
  \toprule
  $(\boldsymbol{\alpha}, \boldsymbol{\omega})$ & \textbf{Distance effective} & \textbf{En clair} \\
  \midrule
  $(*, 0)$ & $\tilde{D}_s$ & Que du vectoriel (courbes ignorées) \\
  $(0, 1)$ & $\tilde{D}_0$ & Que le niveau des courbes \\
  $(1, 1)$ & $\tilde{D}_1$ & Que la forme des courbes \\
  $(0.5, 0.5)$ & Mix de tout & Compromis global \\
  $(0, 0.3)$ & $\sqrt{0.3 \cdot \tilde{D}_0^2 + 0.7 \cdot \tilde{D}_s^2}$ & Surtout vectoriel, un peu de niveau \\
  \bottomrule
\end{tabular}
\end{center}

\subsection{Choix des paramètres : grid search 2D}

On explore toutes les combinaisons
$(\alpha, \omega) \in \{0, 0.1, \ldots, 1\}^2$ ($= 121$ points). Pour chacune :
\begin{enumerate}[nosep]
  \item Calculer $\Dp(\alpha)$.
  \item Calculer $\Dw(\alpha, \omega)$.
  \item Appliquer PAM ($k$ clusters) et mesurer la \textbf{silhouette moyenne}.
\end{enumerate}
On retient le couple $(\alpha^*, \omega^*)$ qui \textbf{maximise la silhouette}.

\begin{remarque}[Critère de sélection]
On utilise la silhouette (critère intrinsèque, pas de vérité terrain) pour
choisir les paramètres. L'ARI (qui nécessite les vrais labels) n'est calculé
qu'a posteriori pour évaluation. Cela garantit que la méthode reste non
supervisée.
\end{remarque}

\subsection{Propriétés}

\begin{itemize}[nosep]
  \item $\Dw$ est une \textbf{vraie distance} (mêmes arguments que $\Dp$).
  \item \textbf{Cas dégénérés} : si $\omega^* = 0$, le fonctionnel est inutile
    sur ce jeu de données. Si $\omega^* = 1$, le vectoriel est inutile.
  \item Le grid search est exact sur la grille mais ne teste pas les valeurs
    entre les pas de 0.1. Affiner la grille (pas de 0.05) est possible mais
    coûteux.
\end{itemize}

% ────────────────────────────────────────────────────────────────────────
\section{$\DK(\alpha)$ : distance à noyaux (stratégie C)}
% ────────────────────────────────────────────────────────────────────────

\subsection{Idée}

Au lieu de pondérer des distances, on passe par des
\textbf{fonctions noyaux} (kernel functions) qui mesurent la
\emph{similarité} (pas la dissimilarité). On les combine par produit, puis on
reconvertit en distance.

\subsection{Noyaux gaussiens}

\begin{definition}[Noyau gaussien]
Pour une distance $D$ et une largeur de bande $\sigma > 0$ :
\[
  K(i,j) = \exp\!\left( -\frac{D(i,j)^2}{2\sigma^2} \right).
\]
\end{definition}

\begin{itemize}[nosep]
  \item $K(i,j) \in [0, 1]$ : 1 si identiques, $\to 0$ si très éloignés.
  \item $\sigma$ contrôle la \textbf{portée} : petit $\sigma$ = seuls les très
    proches ont $K \approx 1$ ; grand $\sigma$ = tout le monde est similaire.
\end{itemize}

\subsection{Construction de $\DK$}

\begin{definition}[$\DK(\alpha)$]
\begin{align}
  K_f(i,j) &= \exp\!\left( -\frac{\Dp(\alpha)_{ij}^2}{2\sigma_f^2} \right),
  & \sigma_f &= \mathrm{m\acute{e}diane}\{\Dp(\alpha)_{ij} : i < j\},
  \label{eq:Kf} \\
  K_s(i,j) &= \exp\!\left( -\frac{D_s(i,j)^2}{2\sigma_s^2} \right),
  & \sigma_s &= \mathrm{m\acute{e}diane}\{D_s(i,j) : i < j\},
  \label{eq:Ks} \\
  K(i,j) &= K_f(i,j) \cdot K_s(i,j), & & \text{(produit de noyaux)}
  \label{eq:Kprod} \\
  \DK(i,j) &= \sqrt{K(i,i) + K(j,j) - 2K(i,j)}. & &
  \label{eq:DK}
\end{align}
\end{definition}

\subsection{Pourquoi le produit~?}

Le noyau produit $K = K_f \cdot K_s$ signifie que deux individus ne sont
similaires \textbf{que s'ils le sont sur les deux sources simultanément}.
Si $K_f(i,j) \approx 0$ (courbes très différentes), alors
$K(i,j) \approx 0$ même si $K_s(i,j) = 1$ (vecteurs identiques).

C'est une forme de \textbf{ET logique flou} entre les deux sources.

\subsection{Heuristique de la médiane pour $\sigma$}

Le choix de $\sigma$ est crucial. Trop petit : $K \approx I_n$ (matrice
identité, tout est dissimilaire). Trop grand : $K \approx \mathbf{1}$
(tout est similaire).

L'\textbf{heuristique de la médiane} (Jaakkola \& Haussler, 1999) fixe
$\sigma$ à la médiane de toutes les distances non nulles. Cela garantit
qu'environ la moitié des paires ont $K > 0.5$ et l'autre moitié $K < 0.5$
— un bon équilibre.

\subsection{Paramètre}

$\DK$ n'a \textbf{pas de paramètre $\omega$} : le produit de noyaux combine
les sources automatiquement. Le seul paramètre libre est $\alpha$ (dans $\Dp$),
optimisé par grid search 1D sur $\{0, 0.1, \ldots, 1\}$.

\subsection{$\DK$ est-elle une distance~?}

$\DK$ vérifie positivité, symétrie et inégalité triangulaire si et seulement
si la matrice $K$ est \textbf{semi-définie positive} (s.d.p.). Le noyau
gaussien est s.d.p. Le produit de noyaux s.d.p.\ est s.d.p. (théorème de
Schur). Donc $K = K_f \cdot K_s$ est s.d.p., et $\DK$ est une distance valide.

\newpage
% ============================================================================
% ============================================================================
\part{Comparaison des stratégies}
% ============================================================================
% ============================================================================

\section{Tableau récapitulatif}

\begin{table}[H]
\centering
\resizebox{\textwidth}{!}{%
\renewcommand{\arraystretch}{1.4}
\begin{tabular}{p{2.8cm}ccp{4.5cm}c}
  \toprule
  \textbf{Distance} & \textbf{Source} & \textbf{Paramètres} & \textbf{Méthode de choix} & \textbf{Clustering} \\
  \midrule
  $D_0$ & Fonct. & 0 & — & PAM \\
  $D_1$ & Fonct. & 0 & — & PAM \\
  $D_s$ & Vect. & 0 & — & PAM \\
  \midrule
  $\Dp(\alpha)$ & Fonct. & 1 ($\alpha$) & Grid search (sil.) & PAM \\
  \midrule
  $\Dw(\alpha, \omega)$ & Mixte & 2 ($\alpha, \omega$) & Grid search 2D (sil.) & PAM \\
  $\DK(\alpha)$ & Mixte & 1 ($\alpha$) + $\sigma$ auto & Grid search 1D (sil.) & PAM \\
  \midrule
  \multicolumn{5}{l}{\textit{Stratégie A (pas de matrice de distances)}} \\
  $W = [\xi \,|\, Z]$ & Mixte & $K$ (nb scores) & Var.\ cumulée $\geq 95\%$ & k-means \\
  \bottomrule
\end{tabular}%
}
\caption{Récapitulatif de toutes les approches.}
\label{tab:recap}
\end{table}


\section{Quand utiliser quoi~?}

\begin{center}
\renewcommand{\arraystretch}{1.3}
\begin{tabular}{p{4.5cm}p{8cm}}
  \toprule
  \textbf{Situation} & \textbf{Distance recommandée} \\
  \midrule
  Pas de variables $Z$ & $\Dp(\alpha)$ seul, optimiser $\alpha$ \\
  $Z$ disponible, on veut un contrôle explicite & $\Dw(\alpha, \omega)$ \\
  $Z$ disponible, on veut du <<~sans paramètre~>> & $\DK(\alpha)$ \\
  Peu de données, besoin de rapidité & Stratégie A (k-means) \\
  Les formes comptent plus que les niveaux & Tester $\alpha > 0.5$ \\
  Les niveaux suffisent & $\alpha = 0$ (revient à $D_0$) \\
  \bottomrule
\end{tabular}
\end{center}

\section{Forces et limites de l'architecture}

\paragraph{Forces.}
\begin{itemize}[nosep]
  \item \textbf{Modulaire} : chaque brique peut être remplacée ($D_1 \to D_2$
    pour les dérivées secondes, $D_s$ euclidienne $\to$ Mahalanobis, etc.).
  \item \textbf{Interprétable} : $\alpha$ et $\omega$ ont un sens clair.
  \item \textbf{Toutes les distances sont valides} dans le sens métrique
    (axiomes vérifiés).
  \item \textbf{Cas limites cohérents} : $\omega = 0 \Rightarrow$ vectoriel
    pur, $\alpha = 0 \Rightarrow D_0$, etc.
\end{itemize}

\paragraph{Limites.}
\begin{itemize}[nosep]
  \item Le \textbf{grid search} est coûteux ($121 \times$ PAM pour la
    stratégie B). Alternatives : Bayesian optimization, descente de gradient
    (mais la silhouette n'est pas dérivable).
  \item La \textbf{normalisation par le maximum} est sensible aux outliers.
  \item Le critère de la \textbf{silhouette} peut ne pas correspondre au vrai
    objectif de clustering.
  \item Si $\Dp$ est redondant avec $D_s$ (comme sur Canadian Weather), le
    système ne peut rien y faire — c'est une propriété des données, pas de
    la méthode.
\end{itemize}

\newpage
% ============================================================================
% ============================================================================
\part{Annexe : résultats sur Canadian Weather}
% ============================================================================
% ============================================================================

\section{Valeurs numériques des distances}

\begin{center}
\begin{tabular}{lrrl}
  \toprule
  \textbf{Distance} & \textbf{Min} (hors 0) & \textbf{Max} & \textbf{Unité} \\
  \midrule
  $D_0$ & 3.9  & 532.1 & °C$\cdot$jours$^{1/2}$ \\
  $D_1$ & 0.51 & 6.61  & °C$\cdot$jours$^{-1/2}$ \\
  $D_s$ & 0.07 & 5.68  & sans unité (standardisé) \\
  \bottomrule
\end{tabular}
\end{center}

\section{Résultats du grid search}

\begin{center}
\begin{tabular}{lccccc}
  \toprule
  \textbf{Stratégie} & $\boldsymbol{\alpha^*}$ & $\boldsymbol{\omega^*}$ & \textbf{Sil.} & \textbf{ARI} & \textbf{Interprétation} \\
  \midrule
  $D_0$ seul   & — & — & 0.285 & 0.229 & Fonctionnel seul insuffisant \\
  $D_s$ seul   & — & — & 0.448 & 0.616 & Vectoriel fort \\
  A (FPCA+Z)   & — & — & 0.395 & \textbf{0.748} & Meilleur ARI \\
  B $(\Dw)$    & 0.0 & 0.0 & \textbf{0.448} & 0.616 & $= D_s$ seul \\
  C $(\DK)$    & 0.0 & — & 0.329 & 0.686 & Noyaux ajoutent un peu \\
  \bottomrule
\end{tabular}
\end{center}

\paragraph{Diagnostic.}
$\alpha^* = 0$ et $\omega^* = 0$ : les deux paramètres tombent sur les
\textbf{coins} de la grille. Cela signifie que sur Canadian Weather, les
variables vectorielles (latitude, longitude) sont quasi-suffisantes et les
courbes n'apportent pas d'information complémentaire (voir rapport principal
pour le diagnostic complet).

\end{document}
